\subsection{Ameaças à Validade}

A validade de um estudo trata da confiabilidade dos seus resultados e a medida com a qual estes são reproduzíveis \cite{wohlin_experimentation_2012}. \citeonline{yin_case_study_2009} retrata de quatro aspectos para se buscar maximizar durante um estudo de caso, são eles a validade do constructo, validade interna, externa e confiabilidade.

A validade do constructo, diz respeito ao grau com que as métricas coletadas para os estudo refletem a intenção de se responder à questão de pesquisa \cite{wohlin_experimentation_2012}. Este aspecto foi tratado através da formalização estruturada das métricas através do GQM (descrito na Subseção \ref{subsec:estudo-questao}) e da escolha de modelos de referência para a escolha das métricas à serem coletadas. Desta forma, garantindo a coleta de dados diretamente relacionados às questões de pesquisa.

Uma possível ameaça a validade do experimento realizado durante o estudo foi a escolha das métricas a partir do modelo de referencia \citeonline{nbr9241}, já que foram observadas apenas duas características de qualidade em uso. Uma avaliação mais assertiva poderia buscar coletar métricas que caracterizassem outros aspectos desta área da qualidade, algo que não foi possível neste estudo dado o contexto do produto sendo estudado.

Já a validade externa, trata-se da extensão com a qual os resultados podem ser generalizados e podem interessar individuos fora do contexto especifico do estudo. Para ativingir este objetivo \citeonline{yin_case_study_2009} sugere a replicação em diversos estudos, o que não é possível em virtude do contexto do trabalho único que é esta monografia.

Desta forma, não se conseguiu atingir a validade externa por completo, apenas sendo possível relatar com o máximo de detalhes contexto específico do projeto, como a escala do produto, os processos de desenvolvimento, entre outros aspectos do caso.

Por fim temos a confiabilidade, que trata da medida com a qual os dados e as analises são influenciadas pelos pesquisadores, já que é esperado que outros pesquisadores possam reproduzir os passos relatados e chegar aos mesmos resultados \cite{wohlin_experimentation_2012}. Para se alcançar esse aspecto da qualidade, foram disponibilizados todos os dados coletados e ferramental referente às análises realizadas, além do relado das mesmas nas Subseções \ref{subsec:pre-estudo} e \ref{subsec:experimento}, de fomar a disponibilizar a reprodutibilidade do estudo em outros contextos \footnote{Endereço do repositório: \href{https://github.com/senaarth/tcc-notebooks/blob/main/experiment.ipynb}{https://github.com/senaarth/tcc-notebooks} }.

Um aspecto não pertinente é a validade interna, já que a mesma faz jus a estudos de investigação de relações causais e, como a atual pesquisa se trata de um estudo de caso, este aspecto não se aplica no atual contexto \cite{yin_case_study_2009}. 